\chapter{Radical Cystectomy}
\section*{Overview}
Radical cystectomy removes the urinary bladder and creates a new route for urine to leave the body.
\section*{Common Instruments and Equipment}
The exact setup varies by technique and setting. Common items may include:
\begin{itemize}[nosep,leftmargin=*]
  \item open or minimally invasive pelvic surgery instrument set
  \item bowel stapling and anastomosis instruments for urinary diversion
  \item ureteral stents and ostomy supplies when applicable
\end{itemize}

\section*{Why It Is Done}
It is a major treatment for muscle-invasive or high-risk bladder cancer and selected locally advanced disease.
\section*{How It Is Performed}
The bladder and usually nearby reproductive structures and lymph nodes are removed. The surgeon creates an ileal conduit, continent reservoir, or neobladder from intestine.
\section*{Contraindications}
No single contraindication applies to every patient; the treating team assesses whether the procedure is suitable for the individual and the reason for treatment.

\section*{Precautions}
The care team reviews relevant health conditions, medicines, allergies, and procedure-specific risks, and explains any changes needed for the individual.

\section*{Risks and Recovery}
Bleeding, infection, bowel complications, urinary leak, electrolyte problems, clots, kidney impairment, and long-term stoma or diversion issues are possible.
\section*{References}
\begin{enumerate}\item National Cancer Institute. Bladder Cancer Surgery. 2025.\item Current Medical Diagnosis and Treatment. McGraw-Hill; 2025.\end{enumerate}
\clearpage
